\newpage
\chapter{KESIMPULAN DAN SARAN} \label{Bab V}

\section{Kesimpulan} \label{V.Kesimpulan}
Berdasarkan hasil penelitian yang telah dilakukan, penelitian ini berhasil merancang, mengimplementasikan, dan mengevaluasi model klasifikasi video \textit{deepfake} menggunakan dua pendekatan, yaitu pendekatan spasial berbasis Vision Transformer (ViT) dan pendekatan spasiotemporal berbasis VideoMAE. Evaluasi dilakukan pada \textit{dataset} \textit{WildDeepfake} sebagai \textit{dataset} utama, serta \textit{Celeb-DF v2} sebagai \textit{dataset} eksternal untuk evaluasi lintas-\textit{dataset} secara \textit{zero-shot}. Adapun kesimpulan yang diperoleh adalah sebagai berikut. \par

\begin{enumerate}[noitemsep]
	\item Model klasifikasi video \textit{deepfake} berhasil diimplementasikan pada tingkat klip untuk kedua pendekatan. Pada model spasial, setiap \textit{frame} diproses sebagai citra tunggal menggunakan ViT, kemudian hasil prediksi \textit{frame} diagregasi menjadi prediksi klip. Sementara itu, model spasiotemporal menggunakan \textit{VideoMAE} untuk memproses \textit{frame} dalam satu klip secara langsung, sehingga hubungan spasial dan temporal dapat dipelajari secara bersamaan.
	\item Pada evaluasi internal menggunakan \textit{WildDeepfake}, pendekatan spasiotemporal menunjukkan performa yang lebih baik dibandingkan pendekatan spasial. Pada eksperimen \textit{baseline}, BASE-ST memperoleh \textit{accuracy} sebesar 85,8\%, \textit{F1-score} sebesar 84,5\%, dan \textit{AUC} sebesar 92,3\%, sedangkan BASE-SP memperoleh \textit{accuracy} sebesar 79,2\%, \textit{F1-score} sebesar 76,7\%, dan \textit{AUC} sebesar 86,8\%. Hasil ini menunjukkan bahwa informasi temporal antar-\textit{frame} berkontribusi terhadap peningkatan performa klasifikasi video \textit{deepfake} pada \textit{dataset in-the-wild}.
	\item Hasil pengujian konfigurasi pelatihan serta konfigurasi input spasial dan temporal menunjukkan bahwa performa model dipengaruhi oleh \textit{learning rate}, \textit{loss function}, augmentasi, \textit{dropout}, \textit{scheduler}, resolusi, jumlah \textit{frame}, dan \textit{stride}. \textit{Learning rate} yang terlalu besar dan augmentasi yang terlalu berat cenderung menurunkan performa model. Selain itu, penurunan resolusi membatasi detail visual yang dapat dipelajari model, sedangkan perubahan jumlah \textit{frame} dan \textit{stride} memengaruhi jumlah informasi temporal yang tersedia. Secara umum, \textit{VideoMAE} tetap lebih unggul pada sebagian besar eksperimen internal \textit{WildDeepfake}, terutama karena mampu memanfaatkan hubungan antar-\textit{frame}.
	\item Evaluasi lintas-\textit{dataset} pada \textit{Celeb-DF v2} menunjukkan pola yang berbeda dari evaluasi internal. Pada skenario \textit{zero-shot}, E6-SP memperoleh \textit{accuracy} sebesar 72,9\%, \textit{F1-score} sebesar 68,7\%, dan \textit{AUC} sebesar 77,7\%, sedangkan E6-ST memperoleh \textit{accuracy} sebesar 64,2\%, \textit{F1-score} sebesar 63,7\%, dan \textit{AUC} sebesar 72,6\%. Hasil ini menunjukkan bahwa pendekatan spasiotemporal yang unggul pada \textit{WildDeepfake} tidak selalu menghasilkan generalisasi terbaik pada \textit{dataset} eksternal. Dengan demikian, pendekatan spasiotemporal lebih efektif pada evaluasi internal \textit{WildDeepfake}, sedangkan pendekatan spasial menunjukkan ketahanan yang lebih baik pada evaluasi lintas-\textit{dataset} terhadap \textit{Celeb-DF v2}.
\end{enumerate}

Secara keseluruhan, penelitian ini menunjukkan bahwa pemanfaatan informasi temporal melalui \textit{VideoMAE} memberikan keuntungan pada klasifikasi video \textit{deepfake} ketika data latih dan data uji berada pada distribusi yang relatif sejalan. Namun, pada skenario lintas-\textit{dataset}, fitur spasial \textit{frame-level} yang dipelajari ViT dapat menjadi lebih stabil terhadap perubahan domain. Oleh karena itu, efektivitas pendekatan spasial dan spasiotemporal perlu dilihat berdasarkan konteks evaluasi, yaitu evaluasi internal atau generalisasi terhadap \textit{dataset} eksternal. \par

\section{Saran} \label{V.Saran}
Berdasarkan penelitian yang telah dilakukan, terdapat beberapa saran yang dapat dipertimbangkan untuk pengembangan penelitian selanjutnya. \par

\begin{enumerate}[noitemsep]
	\item Penelitian selanjutnya disarankan untuk melakukan evaluasi lintas-\textit{dataset} pada lebih dari satu \textit{dataset} eksternal. Pada penelitian ini, evaluasi lintas-\textit{dataset} hanya dilakukan pada \textit{Celeb-DF v2} secara \textit{zero-shot}. Penggunaan \textit{dataset} eksternal lain, seperti \textit{FaceForensics++}, \textit{DFDC}, atau \textit{dataset deepfake in-the-wild} lain, dapat memberikan gambaran yang lebih luas mengenai kemampuan generalisasi model spasial dan spasiotemporal pada berbagai karakteristik data.
	\item Penelitian selanjutnya dapat menguji pengaruh degradasi video secara lebih terkontrol. Pada penelitian ini, variasi input difokuskan pada resolusi, jumlah \textit{frame}, dan \textit{stride}. Oleh karena itu, faktor lain seperti kompresi video, \textit{blur}, \textit{noise}, perubahan \textit{frame rate}, dan \textit{post-processing} dapat diuji sebagai eksperimen terpisah agar pengaruh masing-masing faktor terhadap pendekatan spasial dan spasiotemporal dapat dianalisis secara lebih jelas.
	\item Penelitian selanjutnya disarankan untuk mengeksplorasi strategi \textit{sampling} temporal yang lebih beragam. Penelitian ini menggunakan pembentukan klip berpusat di tengah dan \textit{strided temporal sampling} sebagai strategi pengambilan \textit{frame}. Ke depan, strategi lain seperti \textit{multi-clip sampling}, \textit{segment-based sampling}, \textit{random sampling}, atau \textit{saliency-based sampling} dapat digunakan untuk mengetahui apakah pemilihan \textit{frame} yang lebih adaptif mampu meningkatkan performa dan generalisasi model.
	\item Penelitian selanjutnya dapat mengembangkan pendekatan yang lebih \textit{robust} terhadap pergeseran domain. Hasil evaluasi lintas-\textit{dataset} menunjukkan bahwa model yang baik pada \textit{WildDeepfake} belum tentu menghasilkan performa terbaik pada \textit{Celeb-DF v2}. Oleh karena itu, teknik seperti \textit{domain adaptation}, \textit{domain generalization}, \textit{self-supervised fine-tuning}, atau pelatihan menggunakan gabungan beberapa \textit{dataset} dapat dipertimbangkan untuk meningkatkan ketahanan model terhadap distribusi data eksternal.
	\item Penelitian selanjutnya juga disarankan untuk mengeksplorasi strategi peningkatan performa model, terutama karena beberapa hasil eksperimen masih menunjukkan kecenderungan stagnasi pada nilai di bawah atau sekitar 85\%. Upaya yang dapat dilakukan meliputi penyesuaian strategi \textit{fine-tuning}, penggunaan \textit{hard example mining}, eksplorasi \textit{loss function} yang lebih adaptif, penyeimbangan distribusi data, serta pemanfaatan \textit{pre-training} atau \textit{fine-tuning} tambahan pada \textit{dataset deepfake} yang lebih beragam. Strategi tersebut diperlukan agar model lebih mampu membedakan sampel \textit{real} dan \textit{fake} yang memiliki perbedaan visual maupun temporal yang sangat halus.
\end{enumerate}
